\chapter*{Resumen}
\phantomsection
\addcontentsline{toc}{chapter}{Resumen}

Este trabajo presenta el diseño y la evaluación por simulación de un banco de pruebas de distribución cuántica de claves basado en BB84 con codificación \textit{time-bin} para un enlace de dos nodos. El objetivo es estudiar qué información confiable puede obtenerse antes de construir el banco físico y qué resultados exigen revisar el modelo. Se analizan la tasa de error cuántico de bit (QBER), la tasa tamizada, un indicador ideal de posprocesamiento y estimadores asintóticos con y sin estados señuelo. El indicador ideal no se interpreta como tasa secreta para una fuente coherente débil.

La metodología usa SeQUeNCe 0.8.5 con una corrección local documentada para el error de fase del interferómetro. Se construye un escenario Alice--Bob con canal cuántico, canales clásicos, fuente coherente débil y receptor \textit{time-bin}. Se ejecutan cuatro barridos: distancia, eficiencia y conteos oscuros, visibilidad, y comparación con y sin estados señuelo. Cada punto reúne \PThreeRepetitions{} repeticiones deterministas; QBER se informa con intervalo de Wilson y las tasas con intervalo t del 95\%. La corrida exporta datos agregados y por repetición, semillas, conteos, metadatos y hashes.

El control de consistencia entre tasa simulada y referencia analítica cambia de régimen entre \num[round-mode=places,round-precision=2]{\PThreeTransitionLowKm} y \SI[round-mode=places,round-precision=2]{\PThreeTransitionHighKm}{\kilo\meter}; esta transición diagnostica temporización y contabilidad, no una distancia máxima de QKD. Los contrastes de eficiencia y fondo no resuelven diferencias entre extremos, con \(p=\num[round-mode=places,round-precision=3]{\PThreeEffPValue}\) y \(p=\num[round-mode=places,round-precision=3]{\PThreeDarkPValue}\). La visibilidad produce la respuesta de QBER esperada, aunque sus doce puntos carecen de holgura frente a la referencia de tasa. La estimación con señuelos permanece positiva a \SI{90}{\kilo\meter} bajo un modelo asintótico híbrido, pero activa un tratamiento conservador en parte de las corridas. En consecuencia, el aporte es un procedimiento reproducible para orientar ensayos y detectar límites del modelo. No se certifica seguridad ni factibilidad física del Campus Miguelete.

\textbf{Palabras clave:} distribución cuántica de claves, BB84, time-bin, estados señuelo, SeQUeNCe, fibra óptica, QBER, tasa de clave secreta.

\chapter*{Resumen extendido}
\phantomsection
\addcontentsline{toc}{chapter}{Resumen extendido}

El problema abordado por este trabajo es la distribución segura de claves criptográficas en redes de telecomunicaciones. La criptografía simétrica requiere claves compartidas, pero la distribución de esas claves es un punto crítico cuando los usuarios están separados por enlaces públicos. QKD propone resolver esa etapa usando estados cuánticos y un canal clásico autenticado: si un tercero intenta obtener información midiendo estados no ortogonales, introduce perturbaciones que pueden estimarse como errores.

El proyecto se enfoca en BB84 porque es el protocolo de referencia para estudiar DCC y porque su estructura conecta teoría, simulación y diseño experimental. La elección de \textit{time-bin} responde a su relevancia para fibra óptica: los bits se codifican en ventanas temporales y la base conjugada depende de la interferencia entre modos temprano y tardío. Por eso, el desempeño del enlace queda asociado a atenuación, eficiencia de detección, conteos oscuros, resolución temporal y visibilidad interferométrica.

El aporte específico es construir una primera capa de decisión para el banco. La simulación verifica la respuesta interna a variables físicas, conserva evidencia por repetición y contrasta la tasa tamizada con una referencia externa. Cuando ese contraste pierde holgura, el resultado no se descarta ni se convierte en una prestación: se identifica como condición que requiere auditar temporización, ventanas y conteo de eventos. Esta forma de informar resultados negativos reduce el riesgo de trasladar una salida del simulador a una especificación de hardware sin validación.

\chapter*{Acrónimos}
\phantomsection
\addcontentsline{toc}{chapter}{Acrónimos}

\begin{table}[H]
  \centering
  \begin{tabular}{ll}
    \toprule
    Acrónimo & Significado \\
    \midrule
    APD & Fotodiodo de avalancha \\
    BB84 & Protocolo de Bennett y Brassard de 1984 \\
    DCC & Distribución cuántica de claves \\
    IC & Intervalo de confianza \\
    KMS & Sistema de administración de claves, por \textit{key management system} \\
    OTP & Cifrado de libreta de un solo uso, por \textit{one-time pad} \\
    PFI & Proyecto Final Integrador \\
    PNS & División de número de fotones, por \textit{photon number splitting} \\
    PQC & Criptografía poscuántica, por \textit{post-quantum cryptography} \\
    QBER & Tasa de error cuántico de bit, por \textit{quantum bit error rate} \\
    QKD & Distribución cuántica de claves, por \textit{quantum key distribution} \\
    SKR & Tasa de clave secreta, por \textit{secret key rate} \\
    WCS & Fuente coherente débil, por \textit{weak coherent source} \\
    \bottomrule
  \end{tabular}
  \caption{Acrónimos usados en el trabajo.}
  \label{tab:acronimos}
\end{table}
